\documentclass{article}
\usepackage{graphicx} % Required for inserting images
\usepackage[english]{babel}
\usepackage{amssymb}
\usepackage{amsmath}
\usepackage{amsthm}
\usepackage{mathrsfs}
\usepackage{mathtools}
\usepackage{bm}
\usepackage{xcolor}
\usepackage[]{mdframed}
\usepackage{soul}
\usepackage{algorithm}
\usepackage{algorithmic}

\usepackage[colorlinks=true, allcolors=blue]{hyperref}

\usepackage[letterpaper,top=2cm,bottom=2cm,left=3cm,right=3cm,marginparwidth=1.75cm]{geometry}

\DeclareMathOperator*{\argmax}{arg\,max}
\DeclareMathOperator*{\argmin}{arg\,min}

\newtheorem{proposition}{Proposition}
\newtheorem{definition}{Definition}
\newtheorem{assumption}{Assumption}
\newtheorem{lemma}{Lemma}
\newtheorem{theorem}{Theorem}
\newtheorem{remark}{Remark}
\newtheorem{example}{Example}
\newtheorem{corollary}{Corollary}
\newtheorem{property}{Property}


\title{Conformal Signed Distance Function for Safe and Efficient Motion Planning in Dynamic Environment Comment Section}
\author{}
\date{February 2026}

\begin{document}

{\color{teal}Reviewing what I have written here a month ago, I realized that this formulation lacks some ``context". Should we stick to this formulation, or do we need to consider a ``model" that produces the distance field $D_{t+i}$ from an input $P_{t-H+1:t}$, e.g., point cloud stacked across $H$ frames? In other words, should we view our problem as calibrating a model, or just a spatial function? The latter perspective now seems to have limitations, for instance, consider a single distance map, say $D$, and its estimate $\hat{D}$. While the residual can be naively defined as $R = D - \hat{D}$, one observation is that we can generate \textit{infinitely} many synthetic data pairs from $(D, \hat{D})$. For instance, just take any rigid motion $T: \mathbb{R}^2 \to \mathbb{R}^2$ and consider the pair $(TD, T\hat{D})$. While this copy should be viewed the same as the original one, its residual is $TR$. This observation would lead to a totally different formulation of our formulation. For instance, the score function may need to account for the input itself, which is the stacked point cloud frame, to extract some sort of equivariance.}

{\color{red} 
What I meant was that we are essentially calibrating the error between the model
prediction and the true distance function, treating the prediction model as a
black box, regardless of what information it internally uses. In this view, the
focus is on calibrating the model–reality mismatch itself, rather than
explicitly conditioning on the observation.

That said, I do agree that a fully offline calibration, as written in the
current draft, has a theoretical limitation in that it does not explicitly
account for observation-dependent situations, which is why we are considering
an adaptive version going forward.

However, I personally think that the fully offline version is already
sufficient for the 3D motion planning setting. By not conditioning on the
observation, the calibration effectively aggregates all possible prediction
distributions, leading to a more conservative bound. In 3D, this increased
conservativeness is often acceptable, since there typically remains enough
free and safe space compared to 2D environments.

As for the rigid motion issue, in our current setup all distance fields are
defined on a fixed global coordinate frame obtained via SLAM, so we do not
treat rigidly transformed copies as equivalent samples. }


{\color{teal} 
\noindent
\section{Is our direction reasonable?}
Let us return to our setting. Consider a model that takes a stack of $H$ recent point clouds and predicts the signed distance field after $N$ steps. Formally, this would be defined as a model $h: \mathcal{X} \to \mathcal{Y}$, where
$\mathcal{X}$ is a set of measurements, e.g., $\mathbb{R}^{H \times N \times 3}$ with $N$ denoting the number of points within a single point cloud, 
$\mathcal{Y}$ is a space of some functions. Extending this to embrace stochastic models would be straightforward, but let us assume that $h$ is deterministic. (In our case, this would be the composition of detection, tracking, and trajectory forecasting modules.)

Suppose that for each $\mathbf{x} \in \mathcal{X}$, there exists a label (which is possibly noisy) $\mathbf{y} \in \mathcal{Y}$ so that we can compare the predictions with their ground truths.

\noindent
\textit{Q. How should the errors of the model be defined?}

Our current approach defines the \textit{error} (or \textit{score}) in a \textit{functional} sense:
Let $X$ be the underlying space on which each $\mathbf{y}\in \mathcal{Y}$ is defined. Then, we define
\begin{equation*}
\ell (\mathbf{y}_1, \mathbf{y_2}) \coloneqq \left|\mathbf{y}_1 - \mathbf{y_2} \right|
\end{equation*}
or its \textit{asymmetric} version as we have defined in our project. (The operations are done in a pointwise manner.)
\noindent
While this seems reasonable at a first glance, considering this \textit{functional} error may be vacuous. 
\subsection{How are two errors compared?}
Suppose that we have two samples, $(\mathbf{x}_1, \mathbf{y}_1)$ and $(\mathbf{x}_2, \mathbf{y}_2)$, each of which corresponds to the error $S_1$ and $S_2$, respectively. In other words, we have
\begin{equation*}
S_i = S(\mathbf{x}_i, \mathbf{y}_i) \coloneqq  \ell(h(\mathbf{x}_i), \mathbf{y}_i), \quad i = 1, 2.
\end{equation*}
A straightforward approach (which we have taken in our project) is to rank the scores in a \textit{pointwise} manner:
given $x\in X$, we ask whether $S_1(x) > S_2(x)$ and utilize this rank to analyze  statistics of a model. Unfortunately, such a comparison is a non-sense. The idea makes sense only when the score distribution $S(x)$ over the dataset is distinguishable from that of $S(x')$.}

\color{red}{
To clarify, our approach does not rank error functions in a pointwise manner.
Each sample corresponds to an entire error field
$S_{t+i\mid t}(\cdot)$, which is treated as a single functional observation.
Rather than comparing values $S(x)$ at individual locations,
we construct a functional conformal prediction set
$U_{t+i\mid t}(\cdot)$ such that
\[
\mathbb{P}\!\left[
S_{t+i\mid t}(x) \le U_{t+i\mid t}(x)\ \ \forall x\in\mathcal{X}
\right] \ge 1-\alpha.
\]
To enable conformal calibration, we project each error field onto a finite-dimensional subspace of the underlying Hilbert space and define a scalar conformity score on the resulting coefficient vector.
Importantly, this scalar score serves only as a summary statistic for comparing entire functions,
and the conformal procedure operates at the level of function-valued observations rather than pointwise errors.
}

\color{teal}{
However, note that for any $x$ and $x'$, there always exists a rigid euclidean transformation that sends $x$ to $x'$. Therefore, if $h$ is equivariant under the action fo $\mathrm{E}(n)$ ($n$-dimensional Euclidean group), then the distribution of $S(x)$ would be identical for all $x \in X$. To make this rigorous, let $\mathcal{D}$ be the data distribution (i.e., a probability distribution over $\mathcal{Z}\coloneqq \mathcal{X} \times \mathcal{Y} $). This induces a functional score distribution over $\mathcal{Y}$, namely the push-forward measure of  $\mathcal{D}$ through $S$: $\mathcal{S} = S_{\#} \mathcal{D}$. Note that $\mathcal{S}$ is a distribution of functions defined over $X$. Now fix $x \in X$ and define the evaluation $\pi_x: \mathcal{Y} \to\mathbb{R}$ as
\begin{equation*}
\pi_x(\mathbf{s}) = \mathbf{s}(x),
\end{equation*}
and define $\mathcal{S}_x \coloneqq (\pi_x)_{\#} \mathcal{S} $. We can show that $\mathcal{S}_x$ is identical for all $x$ using the equivariance of the dataset and the predictor $h$:
\begin{equation*}
h(\rho \cdot \mathbf{x}) = \rho\cdot \mathbf{y},\quad \rho_{\#}\mathcal{D} = \mathcal{D}, \quad \rho \in \mathrm{E}(n)
\end{equation*}
\begin{proposition}
$\mathcal{S}_x = \mathcal{S}_{x'}$ for all $x, x'\in X$
\end{proposition}
\begin{proof}
Given $x, x'\in X$, choose $\rho \in \mathrm{E}(n)$ s.t. $\rho(x) = x'$ (A translation would suffice). Then, we have
\begin{flalign*}
\mathcal{S}_{x'} &= (\pi_{x'})_{\#}\mathcal{S} \\
&= (\pi_{x'})_{\#}S_{\#}\mathcal{D} \\
&= (\pi_{x'})_{\#}S_{\#}\rho_{\#}\mathcal{D} & (\text{equivariance of }\mathcal{D})\\
&= (\pi_{x'})_{\#}(S\circ \rho)_{\#}\mathcal{D}  \\
&= (\pi_{x'})_{\#} (\rho \circ S)_{\#}\mathcal{D} & (\text{equivariance of }S) \\
&= (\pi_{x'} \circ \rho)_{\#} S_{\#}\mathcal{D} \\
&= (\pi_x)_{\#} S_{\#} \mathcal{D} & (\rho(x) = x') \\
&= \mathcal{S}_{x}
\end{flalign*}
\end{proof}

The proposition implies that our score definition is not inherently functional; rather its distribution collapses to a distribution over $\mathbb{R}$, merely providing a massive amount of redundancy.

\subsection{Alternative Score Design?}
It may be natural to assume that the data distribution is far from equivariant under $\mathrm{E}(n)$. For instance, lidars are less likely to accurately capture surfaces that are far from them, leading to sparsity. This often makes the data distribution uneven along translation. This also applies to the case where we have perfect obstacle states instead of raw point clouds.
In this context, a possible score candidate would be the set function
\begin{equation*}
\ell(\mathbf{y}_1, \mathbf{y}_2) = \{ \left| R\cdot \mathbf{y}_1 - R\cdot \mathbf{y}_2 \right|: R\in \mathrm{O}(n) \} 
\end{equation*}
which still accounts for the rotation invariance.
Another scenario occurs in online settings: the dataset $\mathcal{D}_{\text{online}}$ is temporally consistent, resulting in high spatial consistency. (That is, the score distribution would be different across $x$, and also they would be more likely to be consistent across time.)
}

\end{document}